\hypertarget{appendix-adfd6397af8d}{}
\section*{Relocation versus move}
\chapter{Relocation versus move}\label{appendix-adfd6397af8d-relocation-versus-move}

\textbf{Curriculum depth:} Introduced · \textbf{First seen:} C009 · \textbf{Latest development:} C009

\section{Definition and mechanism}\label{appendix-adfd6397af8d-definition-and-mechanism}

A move leaves the source object alive; relocation ends its source lifetime as part of transfer. Move-construction plus old-object destruction is a relocation pattern, not general permission to memmove nontrivial objects. Proposal/toolchain facilities are not assumed portable in C++23.

\section{Prerequisites}\label{appendix-adfd6397af8d-prerequisites}

\begin{itemize}
\tightlist
\item
  \hyperref[appendix-5636651017c2]{Move operations and moved-from states} --- Traits, failure policy and relocation must distinguish source survival from source destruction.
\item
  \hyperref[appendix-a82f9c438857]{Object construction and storage reuse} --- A new destination requires a valid lifetime, independently of available bytes.
\end{itemize}

\section{Used by}\label{appendix-adfd6397af8d-used-by}

Later sources will extend this node when their guarantees are developed.

\section{Recall and next depth}\label{appendix-adfd6397af8d-recall-and-next-depth}

Identify how a self-pointer could break representation relocation. Pending: versioned standard facilities and actual toolchain support. Written coverage does not demonstrate personal mastery.

\section{Source and navigation}\label{appendix-adfd6397af8d-source-and-navigation}

\href{Continuous\%20Notes/Computing\%20Foundations\%20-\%20From\%20Bits\%20to\%20C++\%20and\%20Systems\%20APIs\#From\%20three\%20repairs\%20to\%20five\%20decisions\%20to\%20zero\%20boilerplate}{Teaching mechanism and examples} · \hyperref[appendix-fd090d982a15-m009---parts-127-147]{Exact coverage} · \href{../02-Foundations-II/main.pdf\#nameddest=appendix-4ad47bdd269a}{Dependency Map (02-Foundations-II, appendix-4ad47bdd269a)}
